\begin{problem}[第30届全国中学生物理竞赛决赛][太空物资输送与引力离心力做功]{11}
太空中有一飞行器靠其自身动力维持在地球赤道的正上方 $L = \alpha R_{e}$ 处，相对于赤道上的一地面物资供应站保持静止。这里， $R_{e}$ 为地球的半径， $\alpha$ 为常数， $\alpha >\alpha_{m}$ ，而
\begin{equation*}
\alpha_ {m} = \left(\frac {G M _ {e}}{\omega_ {e} ^ {2} R _ {e} ^ {3}}\right) ^ {1 / 3} - 1,
\end{equation*}
$M_{e}$ 和 $\omega_{e}$ 分别为地球的质量和自转角速度， $G$ 为引力常数。设想从供应站到飞行器有一根用于运送物资的刚性、管壁匀质、质量为 $m_{p}$ 的竖直输送管，输送管下端固定在地面上，并设法保持输送管与地面始终垂直。推送物资时，把物资放进输送管下端内的平底托盘上，沿管壁向上推进，并保持托盘运行速度不致过大。忽略托盘与管壁之间的摩擦力，考虑地球的自转，但不考虑地球的公转。设某次所推送物资和托盘的总质量为 $m$。
\begin{subproblem}
在把物资从供应站送到飞行器的过程中，地球引力和惯性离心力做的功分别是多少？
\end{subproblem}
\begin{subproblem}
在把物资从供应站送到飞行器的过程中，外推力至少需要做多少正功？
\end{subproblem}
\begin{subproblem}
当飞行器离地面的高度（记为 $L_{0}$ ）为多少时，在把物资送到飞行器的过程中，地球引力和惯性离心力所做功的和为零？
\end{subproblem}
\begin{subproblem}
如果通过适当控制飞行器的动力，使飞行器在不输送物资时对输送管的作用力恒为零，在不输送物资的情况下，计算当飞行器离地面的高度为 $L = \alpha R_{e} (\alpha > \alpha_{m})$ 时，地面供应站对输送管的作用力；并对 $L > L_{0}$ 、 $L = L_{0}$ 、 $\alpha_{m} R_{e} < L < L_{0}$ 三种情形，分别给出供应站对输送管作用力的大小和方向。
\end{subproblem}
\end{problem}